% !TeX root = ../../main.tex
\chapter{实验口径与复现说明}
本附录将三项研究中容易混淆的计量口径集中说明，供后续复现实验和论文修订使用。正文中的性能数值来自所提供论文最终稿及其图表；本附录提出的记录方式与补充对照属于复现方法，不表示已经重新运行全部原始实验。各章的构造示例仅用于机制推导，与实测数据分别标注。

\section{输入数据与参考结果的固定}
\subsection{稀疏检索数据}
检索实验应保存原始矩阵、查询集合与预处理参数之间的对应关系。矩阵维度和非零数能够快速检查输入规模，但不能唯一确定数据；文件校验值、生成脚本版本及随机种子有助于识别实际使用的内容。若包含人工生成矩阵，应同时记录非零位置分布、数值分布和每行长度生成方式，仅给出总体密度不足以复现访问局部性。

原始 FP32 矩阵与剪枝、量化及聚类后的表示应保持稳定行号映射。对最终结果中的每个行号，都应能够回溯到同一原始向量。如果重排改变了物理行顺序，需保存逻辑行号到物理位置的映射，并在结果输出时恢复。否则，数值正确的分数也可能对应错误的对象，造成难以由吞吐测试发现的功能问题。

参考 Top-K 应固定得分函数与并列规则。若原始得分为内积，不能在生成参考时无意加入行范数归一化；若采用余弦相似度，则需要规定零向量处理。并列分数可采用稳定行号作为次级排序键，或者采用允许等价结果集合的质量定义，但参考与被测系统必须一致。对于接近舍入边界的分数，还应记录参考的累加精度及容差规则。

\subsection{数值求解数据}
CG 复现除了矩阵，还需要右端项、初始解、预条件器与停止条件。来自同一矩阵集合不代表求解轨迹相同；右端项生成方式不同就可能改变局部状态分布和迭代次数。矩阵专属策略的训练设置也应与评价设置一并保存，明确哪些信息在训练中已可见。

对称正定是本章 CG 研究的适用前提。矩阵文件中的存储方式可能只包含一个三角部分，也可能包含完整对称结构；加载器是否展开对称项，会影响非零数、格式转换和访问成本。复现时应检查加载后的实际矩阵，而不是仅依赖文件名或数据集标签。若使用预条件器，其构建精度和应用精度也属于输入契约。

\begin{table}[htbp]\centering\caption{输入与质量参考的最小记录集合}
\begin{tabularx}{\textwidth}{p{2.6cm}XX}\toprule
任务&输入记录&参考结果记录\\\midrule
Top-K 检索&矩阵、查询、逻辑行号、预处理版本&得分函数、K、精确参考集合、并列规则\\
聚类选择&上述数据及簇目录、代表和物理布局&探测区域、局部候选及最终输出\\
CG 求解&矩阵、右端项、初值、停止设置&参考解或残差、迭代数、成功状态\\
Jacobi-PCG&上述数据及对角预条件设置&预条件停止量、真实残差、完整耗时\\\bottomrule
\end{tabularx}\end{table}

\section{性能测量与统计汇总}
\subsection{主机与设备的计时范围}
设备事件适合测量内核和设备流水，但可能不包含主机准备、提交和同步等待。墙钟计时更接近完整调用成本，却容易混入一次性初始化。因此，应同时说明计时起止点、同步方式、预热次数以及是否复用已分配缓冲区。对重复查询和迭代求解，应将初始化与稳态执行分开报告。

FPGA 与 GPU 都可能在首次使用时存在加载、缓存或功率状态变化。仅记录第一轮或仅选择最快一轮均可能产生偏差。可按固定预热协议后测量多次，保留原始样本，再报告中位数、均值或几何平均等预先确定的统计量。测量次数和波动范围应与结果一起保存，不能事后按最有利的统计方式选择结论。

在线求解是否包含结束后的真实残差重算，需要特别说明。EvoCG 原论文的 FP64 真实残差检查用于验证最终质量，未计入报告的在线求解时间。若部署系统要求每次在线检查，应另行测量检查开销并给出包含它的总时间；这属于部署口径变化，不能在不补测的情况下修改实验数字。

\subsection{加速比与平均方式}
设第 $i$ 个任务的参考时间为 $t_i^{(b)}$，被测时间为 $t_i^{(s)}$，逐任务加速比为 $s_i=t_i^{(b)}/t_i^{(s)}$。几何平均为
\begin{equation}
 \bar s_{\mathrm{geo}}=\exp\left(\frac1N\sum_{i=1}^N\ln s_i\right).
\end{equation}
它与总时间之比 $\sum_i t_i^{(b)}/\sum_i t_i^{(s)}$含义不同。前者给每个任务的倍率相同权重，后者更受耗时较长任务影响。论文应沿用实验的汇总口径，并在需要部署评估时补充符合实际请求频率的加权结果。

\begin{table}[htbp]\centering\caption{汇总指标与适合解释的对象}
\begin{tabularx}{\textwidth}{p{3cm}XX}\toprule
指标&主要含义&常见误用\\\midrule
逐任务加速比&特定输入下的时间改善&只展示最快案例代表全部输入\\
几何平均加速比&任务集合中的典型倍率&与累计总时间节省混同\\
整体 QPS&给定提交方式下的处理能力&推断低负载单请求延迟\\
P99 延迟&指定负载下的高分位时间&缺少到达模式仍跨系统直接比较\\
平均 Recall@K&查询集合上的平均覆盖&认为每条查询都达到该阈值\\
成功求解率&满足停止条件的样本比例&把失败上限当作成功时间\\\bottomrule
\end{tabularx}\end{table}

\subsection{能耗与功耗边界}
能效比较应明确功率测量覆盖板卡、处理器还是整机。板卡功耗与整机功耗不能直接混为同一个分母。若使用平均功率乘以执行时间估计能量，应使功率窗口与执行窗口对应，并说明是否包含空闲功率与主机参与部分。短任务尤其容易受到采样间隔限制，需要足够长的重复窗口。

选择性计算降低访问量后，设备功率未必等比例下降，但执行时间缩短仍可能减少能量。相反，索引或训练的一次性成本可能增加总能耗。实验结果表明的板卡功耗用于解释在线实现，长期部署能效仍需结合有效复用次数与更新周期评估，不能只用一次查询的时间乘以额定功率作为最终结论。

\section{各研究章节的收益归因路径}
\subsection{AccelES 的分层对照}
建议首先固定 FP32 输入、查询集合与目标 K，以完整 FP32 计算产生参考。随后依次考察仅量化、量化加 ReSparse 以及增加高精度重算的配置，同时保存预筛候选集合和最终集合。这样可以判断计算量下降是否由数值字段压缩、非零数量减少或精排范围缩小带来。

对于硬件贡献，应在候选预算相同的前提下比较传输格式、计算核心和排序配置。若某配置改变了候选集合，其时间差也包含算法工作量差异，应单独标记。高精度重算的工作量需要按原始行的非零数统计，局部排序损失则需要与数值扰动损失分开识别。

\subsection{ParetoES 的范围与均衡对照}
先固定量化与非零剪枝配置，改变探测预算，得到质量与实际扫描量之间的关系；再固定已选择的簇集合，比较不同布局和负载分配。前一对照解释选择性计算的算法收益，后一对照解释在相同工作量上的执行收益。

热点均衡应同时保存用于画像的查询和未参与画像的测试查询。对两组分别统计核间平均负载、最大核负载和查询时间，可判断映射是否只适配画像样本。若测试查询与画像查询存在分布变化，应记录变化的构造方式，保留原始访问矩阵，以便重新分析相关热点和异常尾部。

\subsection{EvoCG 的精度与执行对照}
先以相同 Tile-CSC、统一内核、融合和图执行构成 FP64 参考，再启用精度策略，记录完整求解时间、迭代数和最终真实残差。这样得到的差异包含精度收益及其控制成本，与跨库比较具有不同含义。进一步改变推理位置、复用间隔或状态特征时，应保持其余条件固定。

对迁移策略，必须保留失败样本。成功子集上的加速可以说明条件性收益，但应同时报告全部样本数量和失败数量。若要比较策略的总体部署价值，需要指定失败后的处理流程并测量完整成本；本初稿没有将尚未实施的回退计入原论文加速结果。

\begin{table}[htbp]\centering\caption{三项研究的收益归因记录}
\begin{tabularx}{\textwidth}{p{2cm}XXX}\toprule
工作&算法层记录&执行层记录&质量层记录\\\midrule
AccelES&量化参数、剪枝、候选集合&扫描字节、重算非零数、排序成本&粗筛覆盖、最终 Recall、重算分数\\
ParetoES&聚类、探测与实际扫描范围&簇跳转、核负载、关键路径&区域覆盖、条件保留、最终 Recall\\
EvoCG&状态、动作与复用轨迹&内核时间、推理、同步与图执行&迭代轨迹、真实残差、失败状态\\\bottomrule
\end{tabularx}\end{table}

\section{论文材料与数据一致性说明}
\subsection{版本选择与图表来源}
AccelES 以所提供的 IEEE 最终 PDF 为主要结果依据，图件使用对应录用版源文件中的矢量 PDF；ParetoES 以 ISCA 最终稿为主要结果依据，并使用配套图件；EvoCG 以 PPoPP 最终稿为依据，从其 PDF 裁切保留矢量图。早期投稿稿件与演示文稿用于理解研究动机和系统结构，不覆盖最终稿中已经调整的实验设置。

演示文稿备注包含口头报告的说明线索，可帮助展开设计动机，但口头简述往往省略质量条件、基线和统计集合。因此，定量结论以论文正文及图表为准。项目中的素材来源清单记录每张复用图的原文件及必要裁切位置，便于后续替换高清原图或改绘中文标注。

\subsection{容易混淆的数值口径}
AccelES 最终稿不同位置存在加速比四舍五入或版本差异，本初稿统一采用最终摘要和实验部分的完整重算配置数值，不混用早期概述中的近似数字。ParetoES 采用最终稿中的质量区间与原始 QPS，避免把配置间提升百分比和绝对吞吐混淆。

EvoCG 的47 矩阵默认求解集合、21 个图中代表矩阵、23 个迁移获益矩阵以及40 个存储非离群矩阵对应不同统计问题。正文分别标明这些集合，不将子集结果扩展为所有矩阵的结果。真实残差验证、训练成本和在线时间也保持分开报告，便于后续补充完整部署评估。

\subsection{后续补测的优先顺序}
优先补测应围绕现有结论中最影响适用边界的问题，而不是重复已充分展示的峰值性能。对前两项研究，应关注未参与校准的查询、接近全召回的工作点和动态数据变化；对 EvoCG，应关注不同右端项、多随机种子、策略迁移失败和包含质量检查的完整求解成本。

这些补测需要对应明确假设。例如“热点均衡在查询分布变化后仍有收益”应由独立变化查询集及重布局成本支持；“矩阵专属策略可以复用于同一矩阵的不同右端项”应由固定策略的跨右端项实验支持。按照假设组织证据，能够使新增实验真正完善适用条件，而不只是增加图表数量。
